%% 01-pendahuluan.tex — bab dokumen contoh publikasi-proposal (hal. 57-58 (proposal), hal. 58-60 (laporan dan repositori)).
\chapter{PENDAHULUAN}
\section{Latar Belakang}
\section{Rumusan Masalah}
Pertanyaan penelitian dituliskan singkat, padat, dan sistematis sebagai dasar penyusunan artikel.

Kajian pustaka disusun dari sumber terkini \citep{creswell2018research, patten2017understanding}
dan dari pedoman resmi universitas \citep{unnes2024panduanakademik}. Regulasi penjaminan mutu
menjadi dasar kebijakan \citep{permendikbudristek2023}.
 Contoh rujukan gaya MLA: \citemla[23]{patten2017understanding}.

\section{Tujuan}
Tujuan penulisan artikel dipastikan terjawab pada abstrak, pembahasan, dan simpulan.

\section{Manfaat}
\section{Kebaruan}
Kebaruan dijelaskan melalui pemetaan area penelitian dan posisi studi ini di dalamnya.
